% =====================================================================
%  Függelék: a módszer és a kiértékelés részletei (a 3. és 4. fejezetből)
% =====================================================================

\chapter{A módszer és a kiértékelés részletei}
\label{app-modszer}

Ez a függelék a 3. és a 4.~fejezet részleteit tartalmazza; a paraméterek
alapértékeit \az{\ref{app-reprodukcio}}.~függelék foglalja össze.

% ---------------------------------------------------------------------
\section{A GPU-memória becslése}
\label{app-vram}
% ---------------------------------------------------------------------

A szövegelőállítás közbeni memóriaigény
\begin{equation}
M \;=\; W \;+\; \mathrm{ctx}\cdot m_{kv} \;+\; M_0,
\label{eq:vram}
\end{equation}
ahol $W$ a kvantált súlyok mérete, $\mathrm{ctx}$ a kontextus hossza tokenben, $M_0$ a
futtatókörnyezet többletigénye, $m_{kv}=2\cdot L\cdot H_{kv}\cdot d_h\cdot 2$ bájt pedig a
KV-gyorsítótár tokenenkénti mérete ($L$ réteg, $H_{kv}$ kulcs--érték fej, $d_h$ elemű vektorok, 2
bájt/elem). A Qwen2.5-7B-re~\cite{qwen25} ($L=28$, $H_{kv}=4$, $d_h=128$) $m_{kv}=56$\,KiB, így a
gyorsítótár 8\,192 tokennél 448\,MiB, 16\,384 tokennél 896\,MiB, 32\,768 tokennél 1,75\,GiB
(\ref{tab:budget}.~táblázat). A PULI-LlumiX 8B~\cite{pulillumix} a Llama-3.1-8B architektúráját örökli. A
Gemma 3 12B rétegeinek többsége csak egy rövid, helyi ablakra (\emph{sliding window}) figyel, és csak
minden hatodik réteg lát a teljes kontextusra, ezért nála a képlet túlbecsül, és a mért értéket
használjuk.

\begin{table}[htbp]
\centering
\caption[Becsült GPU-memóriaigény (8\,192 tokenes kontextus, Q4\_K\_M kvantálás).]{Becsült
GPU-memóriaigény 8\,192 tokenes kontextusnál, Q4\_K\_M kvantálással (kb.\ 4,8 bit/súly) és
$M_0\approx0{,}56$\,GiB többlettel. A súlyok a GGUF-fájlok GiB-ra átszámított mérete, $m_{kv}$ a
publikált architektúrákból számolt érték; a Gemma 3 12B sorában az Ollama által mért foglalás áll.}
\label{tab:budget}
\small
\setlength{\tabcolsep}{4pt}
\begin{tabular}{@{}lccccc@{}}
\toprule
Modell & $L/H_{kv}$ & Súlyok & $m_{kv}$ & Összesen & Kártya \\
\midrule
Qwen2.5-3B  & 36/2 & 1,8\,GiB & 36\,KiB  & 2,6\,GiB  & 8\,GiB \\
Qwen2.5-7B  & 28/4 & 4,4\,GiB & 56\,KiB  & 5,4\,GiB  & 8\,GiB \\
Llama-3.1-8B~\cite{grattafiori2024llama3} (= PULI-LlumiX 8B) & 32/8 & 4,6\,GiB & 128\,KiB & 6,1\,GiB & 8\,GiB \\
Gemma 3 12B~\cite{gemma3} & 48/8 & 7,5\,GiB & ($\le$\,384\,KiB) & $\approx$\,7,8\,GiB (mért) & 12\,GiB \\
Qwen2.5-14B & 48/8 & 8,4\,GiB & 192\,KiB & 10,4\,GiB & 12\,GiB \\
\bottomrule
\end{tabular}
\end{table}

% ---------------------------------------------------------------------
\section{A generálási módszer részletei}
\label{app-bloom}
% ---------------------------------------------------------------------

\paragraph{Bloom-kiosztás.} A nehézségi szinthez tartozó eloszlást (\ref{tab:bloom}.~táblázat)
\emph{legnagyobb maradék módszerrel} (Hamilton-módszer~\cite{balinski2001fair}) alakítjuk egész
darabszámokká: közepes nehézségnél $n=7$ kérdésre a $(2{,}8;\,2{,}8;\,1{,}4)$ szorzatok egészrészei
$(2,2,1)$, a két maradék hely a két legnagyobb törtrészű szinthez kerül, így a kiosztás $(3,3,1)$. A
szinteket véletlenszerűen rendeljük a kérdéshelyekhez, a kész vizsgát a könnyebbtől a nehezebb felé
rendezzük.

\begin{table}[htbp]
\centering
\caption{A nehézségi szintekhez rendelt Bloom-eloszlások és egy tízkérdéses vizsga kiosztása.}
\label{tab:bloom}
\small
\begin{tabular}{@{}lccccc@{\qquad}c@{}}
\toprule
Nehézség & Emlékezés & Megértés & Alkalmazás & Elemzés & Értékelés & Kiosztás ($n=10$) \\
\midrule
Könnyű   & 0,5 & 0,5 & -- & -- & -- & (5, 5, 0, 0, 0) \\
Közepes  & -- & 0,4 & 0,4 & 0,2 & -- & (0, 4, 4, 2, 0) \\
Nehéz    & -- & -- & 0,3 & 0,4 & 0,3 & (0, 0, 3, 4, 3) \\
\bottomrule
\end{tabular}
\end{table}

\paragraph{Darabolás.}\label{app-darabolas} A vágási küszöb a szomszédos mondatok
$\delta_i=1-\cos(\mathbf{v}_i,\mathbf{v}_{i+1})$ távolságainak 85.~percentilise (az \arm{E5b}-ben az
átlag $+\,1{,}5$ szórás), amely a célméret (1\,000 karakter) túllépése után $\tau/(\ell/1\,000)$-re
csökken ($\ell$ az aktuális darab hossza karakterben).

\paragraph{Hibrid keresés.}\label{app-hibrid} Az \arm{E2h} a hat karakterre csonkolt szavakon számolt
BM25-rangsort~\cite{robertson2009bm25} és a sűrű rangsort reciprok rangfúzióval~\cite{cormack2009rrf}
egyesíti: egy darab pontszáma $1/(60+\mathrm{rang}_{\text{sűrű}})+1/(60+\mathrm{rang}_{\text{BM25}})$.

\paragraph{Fogalomgráf.}\label{app-graf} Az \arm{E3} legfeljebb 3\,000 karakteres kötegenként kinyeri a
fogalmakat, tényeket és kapcsolatokat, a 0,92-nél hasonlóbb nevű fogalmakat összevonja, és a fogalmakat
a $\text{kapcsolatok}+\tfrac12\cdot\text{darabok}+\tfrac12\cdot\text{tények}$ centralitás szerint
rendezi; disztraktor-ötletként az azonos szülőjű testvérfogalmakat adja át.

\paragraph{Egy kérdéshely kitöltése.}\label{app-slot} \Az{\ref{alg:slot}}.~algoritmus az \arm{E2} egy
kérdéshelyét tölti ki (vö.\ \ref{sec:verifier}.~szakasz). Ha egy tartalékfogalommal sem sikerül, a
\emph{legjobb} kísérlet (formailag helyes, szivárgásmentes, legkevesebb hibájú) marad meg \emph{nem
ellenőrzöttként} megjelölve.

\begin{algorithm}[htbp]
\caption{Egy kérdéshely kitöltése visszacsatolással (\arm{E2}).}
\label{alg:slot}
\KwIn{kérdéshely $s$ (fogalom, típus, Bloom-szint); elfogadott kérdések $\mathcal{A}$; $k=3$; $r=1$}
\KwOut{kérdés $q$ és az \emph{ellenőrzött} jelző}
$q^{*}\gets\emptyset$\;
\ForEach{fogalom $c \in \bigl(s.\mathit{fogalom},\ \text{az első tartalékfogalom}\bigr)$}{
  $h\gets\emptyset$;\quad $k'\gets k$\;
  \For{$t=0,\dots,r$}{
    $K\gets$ \textsc{Kontextus}$(c,k')$ \tcp*{$\le$ 6\,000 karakter}
    $q\gets$ \textsc{Generál}$(s,c,K,\mathcal{A},h)$ \tcp*{$h$: az előző hibák}
    $h\gets$ \textsc{Ellenőriz}$(q,K)\ \cup$ \textsc{Duplikátum}$(q,\mathcal{A})$\;
    \lIf{$h=\emptyset$}{\KwRet $(q,\ \text{igaz})$}
    $q^{*}\gets$ a jobb $q^{*}$ és $q$ közül\;
    \lIf{a kulcs nem igazolt}{$k'\gets k'+2$}
  }
}
\lIf{$q^{*}$ formája hibás}{$q^{*}\gets$ \textsc{Javít}$(q^{*})$}
\KwRet $(q^{*},\ \text{hamis})$\;
\end{algorithm}

% ---------------------------------------------------------------------
\section{A mérés részletei}
\label{app-adatkeszlet}
% ---------------------------------------------------------------------

\paragraph{Adatkészlet.} Minden referenciakérdéshez rendelkezésre áll a forrásszakasz és a választ
igazoló mondat; az importált forrásokból csak azokat a kérdéseket tartottuk meg, amelyek igazoló mondata
szó szerint szerepel a dokumentumban (\ref{tab:adatkeszlet}.~táblázat).

\begin{table}[htbp]
\centering
\caption[Az \code{eval-v1} adatkészlet és a fő vizsgálatban használt része.]{Az \code{eval-v1} adatkészlet és a fő vizsgálatban használt része. R/K/H: rövid
($\le$\,6\,000), közepes (6\,000--25\,000) és hosszú (25\,000--90\,000 karakteres) dokumentumok.}
\label{tab:adatkeszlet}
\small
\begin{tabular}{@{}llrcrccl@{}}
\toprule
 & & \multicolumn{3}{c}{\code{eval-v1}} & \multicolumn{2}{c}{Fő vizsgálat} & \\
\cmidrule(lr){3-5}\cmidrule(lr){6-7}
Forrás & Nyelv & Dok. & R/K/H & Ref.\ kérdés & Dok. & R/K/H & Licenc \\
\midrule
Saját oktatási szövegek & magyar & 3 & 3/0/0 & 12 & 3 & 3/0/0 & saját \\
OpenStax, EduQG & angol & 12 & 4/4/4 & 74 & 3 & 1/1/1 & CC BY 4.0 \\
Wikipédia, MILQA & magyar & 24 & 8/8/8 & 240 & 24 & 8/8/8 & CC BY-SA 4.0 \\
\midrule
Összesen & & 39 & 15/12/12 & 326 & 30 & 12/9/9 & \\
\bottomrule
\end{tabular}
\end{table}

\paragraph{Oktatói értékelés.}\label{app-oktatoi} A hat közös dokumentumon (3 magyar, 3 angol) futott hat
konfigurációból (Qwen2.5-7B, PULI-LlumiX 8B és Gemma 3 12B; \arm{B-doc-L} és \arm{E2}) futásonként 10
értékelendő és 2 kalibrációs kérdést mintavételezünk (előbb minden dokumentumból egyet, rögzített maggal, csak
a bíró által pontozott kérdések közül). A 72 kérdés minden értékelőnél azonos; több, eltérő kérdéssorrendű lapváltozat
készült, amelyeket összesen öt oktató töltött ki. A lapok vakok: a
konfiguráció neve nem szerepel, az értékelendő kérdések sorrendje lapváltozatonként véletlen, a forrásrészlet
és a rubrika a bíróéval azonos. Az oktatók egymás közötti egyezését Krippendorff-$\alpha$-val~\cite{krippendorff2004}
mérjük (a hiányzó értékeléseket az $\alpha$ kezeli). A bíró és az oktatók egyezését a bíró pontszáma és az
oktatói átlag Spearman-korrelációjával, a bíró pontszáma és a kerekített oktatói medián $\alpha$-jával, a bíró
„alátámasztott” ítéletének és az oktatói többség helyesség $\ge 4$ ítéletének, valamint a két felhasználnám
ítéletnek az egyezését Cohen-$\kappa$-val mérjük. A 95\,\%-os intervallumok a kérdésekre vett bootstrapből
(2\,000 minta) származnak.

\paragraph{Mutatók.}\label{app-mutatok} A \emph{lefedettség}
\begin{equation}
\mathrm{Lef}=\frac{\text{a kérdésekhez (törzs és kulcs) legközelebbi különböző darabok száma}}{K},
\label{eq:cov}
\end{equation}
ahol $K$ a dokumentum darabjainak száma; tíz kérdésnél legfeljebb $10/K$. A további mutatók definíciója:
\begin{itemize}[leftmargin=*,itemsep=1pt,topsep=2pt]
\item \emph{Szivárgás:} a legalább háromszavas kulcs szavainak legalább 80\,\%-a a törzsben is szerepel.
\item \emph{Nem kérdés:} a törzs sem kérdőjelre, sem kettőspontra nem végződik.
\item \emph{Duplikátumráta:} a 0,9-nél nagyobb koszinusz-hasonlóságú kérdéspárok aránya.
\item \emph{Meta-hivatkozás:} az „a szöveg szerint” jellegű fordulatok aránya.
\item \emph{Recall@3:} a referenciakérdéssel futtatott keresés első három találata között szerepel-e az igazoló mondatot tartalmazó darab.
\item \emph{Tematikus hasonlóság:} a referenciakérdések mekkora részéhez van közeli generált kérdés.
\item \emph{Költségmutatók:} a szakaszonkénti futási idő, a hívások és tokenek száma, valamint a GPU-memória alapszintje és csúcsa.
\end{itemize}

\paragraph{Statisztika.}\label{app-statisztika} A hatásnagyság a rangbiszeriális
korreláció~\cite{kerby2014}
\begin{equation}
r_{rb}=\frac{W^{+}-W^{-}}{W^{+}+W^{-}},
\label{eq:rrb}
\end{equation}
ahol $W^{+}$ és $W^{-}$ a pozitív, illetve negatív különbségek rangösszege. $D$ dokumentumnál a
Wilcoxon-próba~\cite{wilcoxon1945} legkisebb kétoldali $p$-értéke $2\cdot 2^{-D}$ (hatnál $0{,}031$,
ötnél $0{,}0625$), Holm-korrekció~\cite{holm1979} után hat összehasonlításnál hatszorosa, $0{,}19$. Két
igen/nem ítéletsor egyezése Cohen-$\kappa$-val~\cite{cohen1960kappa}
\begin{equation}
\kappa=\frac{p_o-p_e}{1-p_e},
\label{eq:kappa}
\end{equation}
ahol $p_o$ a megfigyelt, $p_e$ a véletlenül várható egyezési arány (értelmezés: Landis és
Koch~\cite{landis1977}). A bíró alátámasztottsági ítéletét referenciának tekintve az ellenőrzőre
\begin{equation}
\text{pontosság}=\frac{\text{átengedett és alátámasztott}}{\text{átengedett}},\qquad
\text{felidézés}=\frac{\text{átengedett és alátámasztott}}{\text{alátámasztott}}.
\label{eq:precrec}
\end{equation}
